%% kapa: exam=1 round=30 year=2019 period=2 form=A subject=appraisal_law subject_ko=감정평가관계법규 seq=5 number=5
%% source: Q-Net 감정평가사 기출문제 (cst00309) artlSeq=5205205 file=extracted/1008/5205205/2019년 제30회 감정평가사 1차시험 문제지/2019년 제30회 감정평가사 1차시험 문제지/2019년도 제30회 감정평가사 1차 2교시 A형.hwp
%% conversion: kordoc 4.24.51 -> tools/gen1.py (auto); review=proofread-sonnet; flags=none
\begin{question}{5}
국토의 계획 및 이용에 관한 법령상 도시ㆍ군계획시설 부지의 매수 청구에 관한 설명이다. ( )에 들어갈 내용을 옳게 연결한 것은?
\begin{framedtext}
○ 도시ㆍ군계획시설결정의 고시일부터 ( ㄱ )년 이내에 그 도시ㆍ군계획시설사업이 시행되지 아니하는 경우(실시계획의 인가나 그에 상당하는 절차가 진행된 경우는 제외한다)에 청구할 수 있다. \par
○ 도시ㆍ군계획시설의 부지로 되어 있는 토지 중 ( ㄴ )인 토지(그 토지에 있는 건축물 및 정착물을 포함한다)의 소유자가 청구할 수 있다. \par
\end{framedtext}
\begin{choices}
  \item ㄱ: 10, ㄴ: 지목이 대
  \item ㄱ: 10, ㄴ: 용도지역이 주거지역
  \item ㄱ: 10, ㄴ: 용도지역이 관리지역
  \item ㄱ: 20, ㄴ: 지목이 대
  \item ㄱ: 20, ㄴ: 용도지역이 관리지역
\end{choices}
\end{question}
\answer{1}
